\section{Proposed Method}
\label{sec:method}

Fig.~\ref{fig:architecture} illustrates the proposed \model{} architecture. The design follows three principles. First, local RF distortions should be preserved because device-specific impairments often appear as subtle time-domain and spectral variations. Second, LoRa captures should be modeled as structured RF token sequences rather than as a single global vector. Third, OOB spectral evidence should be injected selectively, since prior LoRa RFFI studies have shown that OOB distortion can be informative, but our ablations indicate that naive OOB concatenation can be unstable. Accordingly, \model{} combines a CNN stem, an RF-HSTU sequence encoder, an auxiliary OOB branch, and an OOB-guided cross-attention fusion module.

\subsection{Input Window and Multi-View RF Representation}

Let a received complex baseband LoRa window be denoted as
\begin{equation}
    \mathbf{x} = [x_1, x_2, \ldots, x_L] \in \mathbb{C}^{L},
\end{equation}
where $L$ is the number of IQ samples in one deterministic evaluation window. In our experiments, $L=8192$. Each window is associated with a device label $y \in \{1,\ldots,C\}$, where $C=24$ for the OSU LoRa dataset.

Instead of using only raw IQ samples, we construct a multi-view RF representation. The main-path views include: 1) raw IQ channels $\Re(\mathbf{x})$ and $\Im(\mathbf{x})$; 2) log-magnitude FFT features; and 3) amplitude--phase features. The OOB view is obtained from spectral bins outside the in-band LoRa bandwidth. These OOB components capture spectral leakage and hardware-related distortion that may not be fully represented by in-band IQ features. This use of OOB evidence follows the observation that transmitter hardware impairments can appear outside the nominal signal band, while our novelty lies in how such evidence is fused with RF token representations.

For cross-day training and ablation experiments, we use per-window z-score normalization for OOB features. For cross-receiver stress tests, we use ratio-based OOB normalization to reduce the effect of receiver gain and energy-scale mismatch. This distinction is important: cross-day experiments evaluate temporal and channel drift under the same receiver setting, whereas cross-receiver experiments intentionally expose receiver front-end calibration mismatch.

Each window is split into $P$ non-overlapping patches:
\begin{equation}
    P = \frac{L}{M},
\end{equation}
where $M$ is the patch size. With $L=8192$ and $M=256$, each window produces $P=32$ RF patches. The patch representation preserves local temporal and spectral structure while keeping the sequence length short enough for efficient edge-side inference.

\subsection{CNN-Stem RF Tokenization}

A lightweight one-dimensional CNN stem is used to extract local RF impairment features from the main-path RF views. This component is motivated by prior CNN-based RF fingerprinting baselines, which are effective at capturing local distortions in IQ or spectrogram-like representations~\cite{jian2020deeplearning_rf_fingerprinting,shen2021lora_spectrogram_cnn}. For the $p$-th RF patch, the CNN stem outputs a $D$-dimensional token:
\begin{equation}
    \mathbf{t}^{\mathrm{main}}_p = f_{\mathrm{stem}}(\mathbf{x}_p),
    \quad p=1,\ldots,P,
\end{equation}
where $\mathbf{x}_p$ denotes the multi-view main-path patch and $f_{\mathrm{stem}}(\cdot)$ is the CNN-stem encoder. Stacking all patch tokens yields
\begin{equation}
    \mathbf{T}_{\mathrm{main}}
    =
    [\mathbf{t}^{\mathrm{main}}_1,\ldots,\mathbf{t}^{\mathrm{main}}_P]
    \in \mathbb{R}^{P \times D}.
\end{equation}

This CNN-stem tokenization differs from a pure CNN classifier in two ways. First, it does not immediately collapse the window into a global representation. Second, it provides token-level features that can be fused with OOB evidence and processed by a sequence encoder.

\subsection{OOB Tokenization}

The OOB branch converts OOB spectral patches into auxiliary tokens:
\begin{equation}
    \mathbf{T}_{\mathrm{oob}}
    =
    [\mathbf{t}^{\mathrm{oob}}_1,\ldots,\mathbf{t}^{\mathrm{oob}}_P]
    \in \mathbb{R}^{P \times D}.
\end{equation}
The OOB tokens are aligned with the main-path tokens at the patch level. This alignment allows the model to compare local RF tokens with corresponding OOB distortion evidence rather than simply appending a global OOB statistic.

A direct concatenation baseline can be written as
\begin{equation}
    \mathbf{T}_{\mathrm{concat}}
    =
    [\mathbf{T}_{\mathrm{main}};\mathbf{T}_{\mathrm{oob}}],
\end{equation}
followed by a projection back to the model dimension. However, our controlled ablations show that this naive fusion is unstable in the small-sample cross-day setting. We therefore use cross-attention to inject OOB information adaptively.

\subsection{OOB-Guided Cross-Attention Fusion}

The core module of \model{} is OOB-guided cross-attention. Main-path RF tokens serve as queries, while OOB tokens provide keys and values:
\begin{align}
    \mathbf{Q} &= \mathbf{T}_{\mathrm{main}}\mathbf{W}_{Q},\\
    \mathbf{K} &= \mathbf{T}_{\mathrm{oob}}\mathbf{W}_{K},\\
    \mathbf{V} &= \mathbf{T}_{\mathrm{oob}}\mathbf{W}_{V},
\end{align}
where $\mathbf{W}_{Q}$, $\mathbf{W}_{K}$, and $\mathbf{W}_{V}$ are learnable projections. The cross-attention output is
\begin{equation}
    \mathbf{A}_{\mathrm{oob}}
    =
    \mathrm{softmax}
    \left(
    \frac{\mathbf{Q}\mathbf{K}^{\top}}{\sqrt{d}}
    \right)
    \mathbf{V},
\end{equation}
where $d$ is the attention head dimension. The fused token sequence is obtained through a residual update:
\begin{equation}
    \mathbf{T}_{\mathrm{fuse}}
    =
    \mathrm{LN}
    \left(
    \mathbf{T}_{\mathrm{main}}
    +
    \mathbf{A}_{\mathrm{oob}}\mathbf{W}_{O}
    \right),
\end{equation}
where $\mathbf{W}_{O}$ is an output projection and $\mathrm{LN}(\cdot)$ denotes layer normalization.

This formulation lets each RF token selectively attend to the OOB evidence that is most relevant to that local patch. Compared with early concatenation, cross-attention avoids forcing the classifier to interpret OOB and main-path features as a single undifferentiated vector. Compared with a fixed gate, it produces a patch-dependent alignment between in-band RF structure and OOB spectral distortion. The fusion design builds on the general attention mechanism~\cite{vaswani2017attention}, but is specialized here for OOB-guided LoRa RFFI.

\subsection{Chirp-Aware Positional Embedding}

LoRa uses chirp spread spectrum modulation, and the relative location of a patch within a chirp-like structure can provide useful inductive bias. We therefore add a lightweight chirp-aware positional embedding to the fused tokens:
\begin{equation}
    \widetilde{\mathbf{T}}_{\mathrm{fuse}}[p]
    =
    \mathbf{T}_{\mathrm{fuse}}[p]
    +
    \mathbf{e}_{\mathrm{chirp}}(c_p)
    +
    \mathbf{e}_{\mathrm{pos}}(r_p),
\end{equation}
where $c_p$ denotes the estimated chirp group index of patch $p$, and $r_p$ denotes its relative intra-chirp position. The embeddings $\mathbf{e}_{\mathrm{chirp}}(\cdot)$ and $\mathbf{e}_{\mathrm{pos}}(\cdot)$ are learned.

This module should be interpreted as an auxiliary LoRa-structure prior rather than the sole source of improvement. As shown later in the fusion/chirp ablation, removing chirp embeddings from the cross-attentive model does not degrade cross-day accuracy, while adding chirp embeddings to naive concatenation does not restore stability. Therefore, the primary mechanism is cross-attentive OOB fusion, with chirp embedding retained as a compatible structural prior.

\subsection{RF-HSTU Sequence Encoder}

The fused and position-enhanced token sequence is passed into an RF-HSTU encoder:
\begin{equation}
    \mathbf{H}
    =
    f_{\mathrm{RF\text{-}HSTU}}
    \left(
    \widetilde{\mathbf{T}}_{\mathrm{fuse}}
    \right)
    \in \mathbb{R}^{P \times D}.
\end{equation}
We use ``RF-HSTU'' to refer to our lightweight RF sequence encoder implemented for patch-level LoRa RFFI. The encoder is designed to model interactions across RF patches while preserving efficient inference. Each block contains token mixing, normalization, and pointwise gating operations over the patch sequence. In contrast to a pure CNN classifier, this encoder allows the model to aggregate local impairment evidence across the full LoRa window before classification.

We do not claim that the sequence encoder alone solves all deployment shifts. In fact, our diagnostic baseline without OOB cross-attention performs below the final model. The purpose of the RF-HSTU encoder is to provide a sequence modeling backbone into which OOB evidence can be injected through cross-attention.

\subsection{Classification Objective and File-Level Voting}

The encoded token sequence is converted into a window-level embedding by mean pooling:
\begin{equation}
    \mathbf{z}
    =
    \frac{1}{P}
    \sum_{p=1}^{P}
    \mathbf{H}[p].
\end{equation}
A linear classifier produces device logits:
\begin{equation}
    \mathbf{s}
    =
    \mathbf{W}_{c}\mathbf{z} + \mathbf{b}_{c},
    \quad
    \mathbf{s}\in\mathbb{R}^{C}.
\end{equation}
The model is trained with cross-entropy loss:
\begin{equation}
    \mathcal{L}
    =
    -\frac{1}{N}
    \sum_{i=1}^{N}
    \log
    \frac{\exp(s_{i,y_i})}
    {\sum_{c=1}^{C}\exp(s_{i,c})},
\end{equation}
where $N$ is the number of training windows in a mini-batch.

Since LoRa RFFI authentication is usually performed at the capture-file level rather than from a single random window, evaluation uses deterministic multi-window voting. For a test file containing $K$ evaluation windows, we average the logits:
\begin{equation}
    \bar{\mathbf{s}}
    =
    \frac{1}{K}
    \sum_{k=1}^{K}
    \mathbf{s}_{k},
\end{equation}
and predict
\begin{equation}
    \hat{y}
    =
    \arg\max_{c}
    \bar{s}_{c}.
\end{equation}
This mean-logits voting protocol reduces window-level randomness and matches the file-level authentication setting used in our main experiments.

\begin{figure}[!t]
    \centering
    \includegraphics[width=\linewidth]{fig1_model_architecture.pdf}
    \caption{Proposed OOB-guided cross-attentive RF-HSTU hybrid architecture. A CNN stem extracts local RF tokens from multi-view LoRa windows. An auxiliary OOB branch produces spectral-distortion tokens. Main-path tokens query OOB tokens through cross-attention, and the fused sequence is processed by an RF-HSTU encoder with chirp-aware positional embeddings. File-level mean-logits voting yields authentication-style device decisions.}
    \label{fig:architecture}
\end{figure}
